% Load variables
\newcommand{\myUni}{Università di Padova}
\newcommand{\myDepartment}{Dipartimento di Matematica ``Tullio Levi-Civita''}
\newcommand{\myFaculty}{Corso di Laurea in Informatica}
\newcommand{\myTitle}{Dall'analisi manuale alla verifica assistita: un sistema basato su Large Language Models per l'accessibilità. ​}
\newcommand{\myDegree}{Tesi di Laurea}
\newcommand{\profTitle}{Prof.ssa}
\newcommand{\myProf}{Ombretta Gaggi}
\newcommand{\graduateTitle}{Laureando}
\newcommand{\myName}{Sofia De Blasi}
\newcommand{\myStudentID}{2111014}
\newcommand{\myAA}{2025-2026}
\newcommand{\myLocation}{Padova}
\newcommand{\myTime}{Dicembre 2026}
% Acronyms
\newacronym{api}{API}{Application Program Interface}
\newacronym{uml}{UML}{Unified Modeling Language}
\newacronym{tsa}{TSA}{Termine solo acronimo}
\newacronym{wcag}{WCAG}{Web Content Accessibility Guidelines}
\newacronym{agid}{AGID}{Agenzia per l'Italia Digitale}
\newacronym{llm}{LLM}{Large Language Model}
\newacronym{eaa}{EAA}{European Accessibility Act}
\newacronym{ict}{ICT}{Information and Communication Technology}

% Glossary
\newglossaryentry{}{
    name={},
    text={},
    sort=,
    description={}
}


% Define custom colors
\definecolor{hyperColor}{HTML}{7D0767}
\definecolor{tableGray}{HTML}{F5F5F7}
\definecolor{veryPeri}{HTML}{6667ab}

% Set line height
\linespread{1.5}

% La classe "book", essendo "twoside", usa di default \flushbottom: per far
% terminare ogni pagina esattamente al margine inferiore, LaTeX allunga gli
% spazi elastici disponibili nella pagina. Poiché "parskip" (v. packages.tex)
% separa i paragrafi con uno spazio elastico anziché fisso, nelle pagine con
% pochi paragrafi questo allungamento diventa molto vistoso (spazi enormi e
% disuguali fra un paragrafo e l'altro). \raggedbottom disattiva questo
% comportamento: il margine inferiore può variare leggermente da pagina a
% pagina, ma lo spazio fra i paragrafi resta uniforme.
\raggedbottom

% Custom hyphenation rules
\hyphenation {
    data-base
    al-go-rithms
    soft-ware
}

% Images path
\graphicspath{{img/}}

% Force page color, as some editors set a grayish color as default
\pagecolor{white}

% Better spacing for footnotes
\setlength{\skip\footins}{5mm}
\setlength{\footnotesep}{5mm}

\setlength{\headheight}{14.5pt}
\setlength{\headsep}{1.2cm} % più margine tra la riga di intestazione e il testo
\addtolength{\topmargin}{-2.45pt}

% Add a subscript G to a glossary entry
\newcommand{\glox}{\textsubscript{\textbf{\textit{G}}}}

% Placeholder visibile per le parti non ancora scritte
\newcommand{\todoplaceholder}[1]{\par\noindent\textcolor{red!65!black}{\textbf{[DA COMPLETARE]} #1}\par}

% Comando di citazione per le fonti giuridiche (leggi, decreti, direttive):
% stampa il campo "usera" del file .bib (tipo di atto, data completa, numero)
% invece della citazione autore-anno/numerica, mantenendo comunque il
% collegamento bibliografico e il backref standard di biblatex.
% Si usa il campo generico "usera" (e non "shorthand") perché "shorthand" è
% un campo riservato di biblatex: se un'entry lo definisce, gli stili
% numerici (incluso "ieee") lo utilizzano al posto del numero progressivo
% come "labelnumber", rompendo la numerazione IEEE e la relativa gabbia di
% impaginazione (hanging indent) della bibliografia. "usera" non ha invece
% alcun significato predefinito per biblatex ed è quindi sicuro da riusare.
% Uso identico a \footcite: \legfootcite[postnote]{chiave-bib}
\DeclareCiteCommand{\legfootcite}[\mkbibfootnote]
  {}
  {\usebibmacro{citeindex}%
   \thefield{usera}%
   \iffieldundef{postnote}
     {}
     {\addcomma\space\printfield{postnote}}}
  {\multicitedelim}
  {\addperiod}

% Ridefinizione di \footcite per le fonti non giuridiche in stile IEEE:
% lo stile "ieee" (biblatex-ieee) definisce di default \footcite senza le
% parentesi quadre che invece usa per \cite nel corpo del testo (assume che
% la nota a piè di pagina renda già evidente il rimando). Qui si vuole che
% SOLO il numero della fonte sia racchiuso tra parentesi quadre, mentre
% l'eventuale postnote (capitolo, paragrafo, pagina) resta fuori dalle
% parentesi, secondo il formato:
%   [3], cap. 4, § 4.2.
% invece di "[3, cap. 4, § 4.2]" o di "3, cap. 4, § 4.2" senza parentesi.
% NOTA DI PORTABILITÀ: "ieee.cbx" (biblatex-ieee) sceglie internamente quale
% variante di stile citazionale caricare per implementare \cite/\footcite
% (tipicamente "numeric-comp", che comprime intervalli di citazioni multiple
% del tipo [3]-[5]); tale variante può differire fra installazioni/versioni
% diverse di biblatex-ieee (ad es. TeX Live vs. MiKTeX), e con essa cambiano
% anche i nomi dei bibmacro interni usati per assemblare la citazione
% (bibmacro come "cite:init"/"cite:comp"/"cite:dump" esistono infatti SOLO
% nella variante "numeric-comp"/"numeric-icomp", non ad esempio in
% "numeric-verb", dove "\footcite" è invece assemblato con un bibmacro
% generico chiamato semplicemente "cite"). Per evitare che questa
% ridefinizione smetta di funzionare a seconda della variante caricata,
% qui non ci si appoggia a NESSUNO di questi bibmacro interni: il numero
% viene stampato direttamente con \printfield{labelnumber} (preceduto dal
% campo, anch'esso opzionale, "labelprefix"), racchiuso in
% \printtext[bibhyperref]{...} per mantenere il collegamento ipertestuale
% e il backref verso la bibliografia. "citeindex", "labelnumber",
% "labelprefix" e il field-format "bibhyperref" sono infatti definiti in
% modo generico in biblatex stesso (non in uno specifico file di stile),
% e sono quindi stabili su qualunque variante/versione di biblatex-ieee.
% Nota: questa versione non comprime in intervalli citazioni multiple nello
% stesso \footcite (es. \footcite{a,b,c}), a differenza della variante
% "numeric-comp": non è un problema per il presente lavoro, dove ogni nota
% a piè di pagina cita di norma una sola fonte per volta.
\DeclareCiteCommand{\footcite}[\mkbibfootnote]
  {\usebibmacro{prenote}}
  {\usebibmacro{citeindex}%
   \bibopenbracket
   \printtext[bibhyperref]{\printfield{labelprefix}\printfield{labelnumber}}
   \bibclosebracket}
  {\multicitedelim}
  {\usebibmacro{postnote}}

% Titolo di capitolo: numero e titolo sulla stessa riga (niente "Capitolo N" a capo)
\titleformat{\chapter}[hang]{\normalfont\Huge\bfseries}{\thechapter}{1em}{}
\titlespacing*{\chapter}{0pt}{50pt}{40pt}

% Improvements the paragraph command
\titleformat{\paragraph}
{\normalfont\normalsize\bfseries}{\theparagraph}{1em}{}
\titlespacing*{\paragraph}
{0pt}{3.25ex plus 1ex minus .2ex}{1.5ex plus .2ex}

% Define use case environment
\newcounter{usecasecounter} % define a counter
\setcounter{usecasecounter}{0} % set the counter to some initial value
% Parameters
% #1: ID
% #2: Nome
\newenvironment{usecase}[2]{
    \renewcommand{\theusecasecounter}{\usecasename #1}  % this is where the display of the counter is overwritten/modified
    \refstepcounter{usecasecounter} % increment counter
    \vspace{2em}
    \par \noindent % start new paragraph
    {\normalsize \textbf{\usecasename #1: #2}} % display the title before the content of the environment is displayed
    \vspace{.5em}
}{
    \medskip
}
\newcommand{\usecasename}{UC}
\newcommand{\usecaseactors}[1]{\textbf{\\Attori Principali:} #1}
\newcommand{\usecasepre}[1]{\textbf{\\Precondizioni:} #1}
\newcommand{\usecasedesc}[1]{\textbf{\\Descrizione:} #1}
\newcommand{\usecasepost}[1]{\textbf{\\Postcondizioni:} #1}
\newcommand{\usecasealt}[1]{\textbf{\\Scenario Alternativo:} #1}

% Define risks environment
\newcounter{riskcounter} % define a counter
\setcounter{riskcounter}{0} % set the counter to some initial value
% Parameters
% #1: Title
\newenvironment{risk}[1]{
    \refstepcounter{riskcounter} % increment counter
    \par \noindent % start new paragraph
    \textbf{\arabic{riskcounter}. #1} % display the title before the content of the environment is displayed
}{
    \par\medskip
}
\newcommand{\riskname}{Rischio}
\newcommand{\riskdescription}[1]{\textbf{\\Descrizione:} #1.}
\newcommand{\risksolution}[1]{\textbf{\\Soluzione:} #1.}

% Apply fancy styling to pages
\pagestyle{fancy}
\fancyhf{}
\renewcommand{\chaptermark}[1]{\markboth{#1}{}} % intestazione in minuscolo (non tutto maiuscolo)
\fancyhead[L]{\leftmark} % Places Chapter N. Chatper title on the top left
\fancyfoot[C]{\thepage} % Page number in the center of the footer

% Adds a blank page while increasing the page number
\newcommand\blankpage{ 
\clearpage
    \begingroup
    \null
    \thispagestyle{empty}
    \hypersetup{pageanchor=false}
    \clearpage
\endgroup
}

% Make glossaries and bibliography
\makeglossaries
% Redefine the format for the glossary entries to be italic
\renewcommand*{\glstextformat}[1]{\textit{#1}\glox}
%\glsaddall

\bibliography{references/bibliography}
\defbibheading{bibliography} {
    \cleardoublepage
    \phantomsection
    \addcontentsline{toc}{chapter}{\bibname}
    \chapter*{\bibname\markboth{\bibname}{\bibname}}
}

% Code blocks handling w/ table of codes
\makeatletter
\ifdefined\NR@chapter
  \expandafter\@firstoftwo
\else
  \expandafter\@secondoftwo
\fi{\patchcmd\NR@chapter}{\patchcmd\@chapter}
  {\addtocontents{lot}{\protect\addvspace{10\p@}}}
  {\addtocontents{lot}{\protect\addvspace{10\p@}}%
   \addtocontents{lol}{\protect\addvspace{10\p@}}}
  {}{}
\makeatother

\renewcommand\listingscaption{Codice}
\renewcommand\listoflistingscaption{Elenco dei codici sorgenti}
\counterwithin*{listing}{chapter}
\renewcommand\thelisting{\thechapter.\arabic{listing}}

% Set up hyperlinks
\hypersetup{
    colorlinks=true,
    linktoc=all, % nell'indice l'intera riga è cliccabile, non solo il numero di pagina
    pdfstartpage=1,
    pdfstartview=,
    breaklinks=true,
    pdfpagemode=UseNone,
    pageanchor=true,
    pdfpagemode=UseOutlines,
    plainpages=false,
    bookmarksnumbered,
    bookmarksopen=true,
    bookmarksopenlevel=1,
    hypertexnames=true,
    pdfhighlight=/O,
    allcolors = hyperColor
}

% I riferimenti a termini di Glossario, Acronimi e abbreviazioni (\gls, \acrshort, ecc.)
% restano cliccabili: il collegamento ipertestuale rimanda alla voce
% corrispondente nell'elenco (Acronimi e abbreviazioni / Glossario) in fondo
% al documento, dove si trova la relativa spiegazione. Questo è il
% comportamento predefinito del pacchetto "glossaries" quando "hyperref" è
% caricato prima di esso (v. config/packages.tex): non è quindi necessaria
% alcuna istruzione aggiuntiva; in precedenza tali link erano disattivati con
% \glsdisablehyper, ora rimosso.

% Set up captions
\captionsetup{
    tableposition=top,
    figureposition=bottom,
    font=small,
    format=hang,
    labelfont=bf
}

% When draft mode is on, the hyperlinks are removed. Useful when printing the document. To enable/disable, uncomment/comment the command
% \hypersetup{draft}

% prevents cleaning up the cache at the end of the run (needed to keep the unused caches, generated by other editions)
% guarded with \@ifundefined so it does not error out on minted/MiKTeX versions
% where \minted@cleancache is not defined (e.g. newer minted releases on MiKTeX)
\makeatletter
\@ifundefined{minted@cleancache}{}{\renewcommand*{\minted@cleancache}{}}
\makeatother

% Break lines in code blocks whe using inputminted
\setminted{breaklines}
% ------------------------------------------------------------------
% Protezione contro caratteri Unicode invisibili incollati per errore
% (tipico quando si copia testo da Word, Google Docs o pagine web: si
% porta dietro caratteri come lo zero-width space, che LaTeX non sa
% interpretare e che causano un errore "Unicode character ... not set
% up for use with LaTeX" apparentemente senza motivo visibile).
% Qui vengono semplicemente resi "no-op" (ignorati), così da non
% bloccare la compilazione se dovessero ricomparire in futuro.
\newunicodechar{​}{}
\newunicodechar{‌}{}
\newunicodechar{‍}{}
\newunicodechar{}{}
\newunicodechar{⁠}{}
% ------------------------------------------------------------------